\section{Related Work}
\label{sec:related}

The toolkit is positioned in the space between five existing strands of
work. We summarize how robopsychology relates to each, and what it
does \emph{not} attempt to do. A more discursive treatment of these
distinctions lives in the repository's \texttt{framework/related-work.md}
document.

\paragraph{Capability benchmarks.}
Benchmarks such as HELM \citep{liang2022holistic}, MMLU
\citep{hendrycks2020measuring}, and BigBench
\citep{srivastava2022beyond} measure model performance at scale across
standardized tasks. They answer \emph{``how does model~A compare to
model~B on task category~$X$?''} Robopsychology does \emph{not}
compare models on tasks. It diagnoses why a specific model produced a
specific unexpected output in a specific context. The two modes are
complementary: a benchmark tells the practitioner that a model scores
some percentage on safety; robopsychology helps them understand
\emph{why it refused their particular request} and \emph{whether that
refusal was justified}.

\paragraph{Red teaming and adversarial evaluation.}
Red teaming \citep{perez2022red, ganguli2022red} deliberately tries to
break models --- eliciting harmful outputs, finding jailbreaks, testing
guardrails. The goal is adversarial: provoke failure under pressure.
Robopsychology is diagnostic, not adversarial. Its prompts explicitly
invite the model to \emph{collaborate in its own diagnosis} rather than
attempting to trick it. Adversarial framing is listed as an
anti-pattern in the method.

\paragraph{Alignment evaluations.}
Alignment evaluations (model cards, system cards, third-party audits)
assess whether a model's behavior broadly conforms to intended values
across many scenarios. They are typically run by the model developer
or an auditing body. Robopsychology is a \emph{user-side} tool. It does
not evaluate whether a model is aligned in general --- it helps
individual users understand what happened in a specific interaction.
It is post-hoc and situational, not pre-deployment and comprehensive.

\paragraph{Behavioral testing frameworks.}
CheckList \citep{ribeiro2020beyond} introduced task-agnostic behavioral
testing for NLP --- testing capabilities, perturbation invariance, and
directional expectations. The idea of testing \emph{behavior} rather
than \emph{accuracy on a benchmark} is foundational to robopsychology's
approach. Where CheckList is offline and batch-oriented (running suites
of templated test cases against a model), robopsychology is online and
interactive: the diagnostician adapts the probe to the specific
behavior observed in a real interaction.

\paragraph{Interpretability.}
Mechanistic interpretability \citep{olah2020zoom, elhage2021mathematical}
seeks to explain model behavior through internal mechanisms ---
attention patterns, neuron activations, circuits. It requires
white-box access. Robopsychology is strictly behavioral and black-box:
the diagnostician has only the model's input--output behavior to work
with. The two are complementary; a mechanistic explanation could in
principle confirm or refute a behavioral diagnosis after the fact.

\paragraph{Sycophancy and behavioral failure modes.}
A growing body of work characterizes specific behavioral failure modes
in language models: sycophancy \citep{sharma2023sycophancy,
perez2022discovering}, hallucination, refusal mismatch (XSTest
\citep{rottger2023xstest}), and reward hacking
\citep{lambert2024rewardbench}. Robopsychology overlaps with this work
in vocabulary but differs in posture: where these papers
\emph{characterize} a failure mode at scale, the diagnostic toolkit
helps a practitioner \emph{determine, on a specific interaction},
whether such a mode is the underlying cause. The toolkit treats
existing failure-mode taxonomies as part of its differential diagnosis
catalog rather than as evaluation targets.

\todo{If room, add a paragraph on clinical-psychology methodology
parallels (structured interview vs.\ psychometric test). Keep at most
2--3 sentences --- the analogy is illustrative, not load-bearing.}
